\clearpage
\fancyhead[L]{\small BOUVARD. JUPITER UND SATURN.}
\SourceLeaf{513}
\begin{center}\small Göttingische gelehrte Anzeigen. No. 177. pp. 1761..1764. 5 November 1808.\\[3pt]\rule{.72\linewidth}{.3pt}\end{center}

Available from Courcier. \emph{Tables astronomiques publiées par le bureau des longitudes de France. Nouvelles tables de Jupiter et de Saturne calculées d'après la théorie de M. }\textsc{Laplace}, \emph{et suivant la division decimale de l'angle droit; par M.} \textsc{Bouvard}. 1808. In quarto. Paris.

These tables must be regarded as a continuation of the solar and lunar tables published in 1806 and already reviewed in issue 95 of these pages for 1807. They are the precious fruit of more than 20 years of work in which \textsc{Laplace} continued to refine his investigations of the perturbations that so complicate the motions of the two largest planets in our solar system. It is well known how successfully this great geometer overcame every difficulty here, as in all the other branches of physical astronomy he made the subject of his research, and made these recalcitrant motions yield to calculation. We owe it to him that we can now calculate this motion almost as accurately as we can observe it. The earlier tables of Jupiter and Saturn, based on \textsc{Laplace}’s first investigations and calculated by \textsc{Delambre}, had even then seldom differed from the sky by as much as half a minute. Since then \textsc{Laplace} has carried the calculation of the perturbations much further, and \textsc{Bouvard}, taking those calculations into account, has determined new elliptic elements from a most careful analysis of all the oppositions observed over half a century. This work has led to agreement three times as close. The latest results of this work, namely the motions of Jupiter and Saturn expressed as formulas, have already been published in part 8 of the Mécanique céleste, p. 337. In the present work we now find them put into tables. We have little to add about their arrangement, since in most respects it agrees with the arrangement already familiar from \textsc{Delambre}’s solar tables. Times are reckoned from midnight instead of noon. The anomalies measured from perihelion and the equations were all made positive by adding a constant quantity to each, and subtracting the sum of all these additions again from the equation of the centre (which was made positive where necessary by adding 400 decimal degrees), the elliptic radius vector and the distance from the north pole used in place of latitude. All these conventions have long met with astronomers’ approval, apart from the new way of reckoning time, which was adopted without any apparent benefit. The use of the decimal division of the quadrant in these 'tables is meant as an attempt to accustom astronomers to it gradually. As long as solar and sine tables with the old division still have to be used alongside them, however, this is more of an inconvenience than a help.

The introduction preceding the tables first describes each of them. It then gives the formulas used to calculate them. These differ from those given in the Mécanique céleste at the place cited only in that all epochs and variable coefficients have been reduced from 1750 to 1800. Finally the introduction gives the full calculation of an observation of Jupiter made on 2 April 1806, in which the tables differ by $20^{\prime\prime}0$ in longitude and by $19^{\prime\prime}7$ in latitude ($6^{\prime\prime}8$ and $6^{\prime\prime}4$ in the sexagesimal division). This is the error in the geocentric position. The introduction then shows how to derive the error in the heliocentric position from it. Since the calculated radius vector has to be regarded as perfectly accurate here, the introduction should have pointed out that the procedure can be appl\SourceLeaf{514}ied with some confidence only where a moderate error in the radius vector has very little effect on longitude and latitude. That means only for distant planets or near opposition, and with moderate inclinations or near the nodes.

Since these tables of Jupiter and Saturn are based entirely on the more recent observations, it is interesting to see how accurately they represent earlier observations. Especially noteworthy in this respect is the conjunction of the two planets observed by \textsc{Ibn Junis} in Cairo on 31 October 1007. According to \textsc{Bouvard}’s calculation, the tables differ from the observation by $11'\,25''$ in longitude and $6'\,56''$ in latitude (sexagesimal division), which scarcely exceeds the possible error of observation. We should also note that a small error in the great equation for both Jupiter and Saturn was found after the present tables appeared. A term had been taken with the wrong sign. This change in the great equation also requires a small change in the epochs and mean motions. The result is that to $q^{\mathrm{IV}}$ (Mécan. cél., vol. 4, p. 338) one must add
\[
\begin{aligned}
&51''98 + t.0''4156\\
&{}- (73''58 - t.0''0103)\sin(5n^{\mathrm{V}}t - 2n^{\mathrm{IV}}t + 5\varepsilon^{\mathrm{V}} - 2\varepsilon^{\mathrm{IV}})\\
&{}+ (47''63 + t.0''0287)\cos(5n^{\mathrm{V}}t - 2n^{\mathrm{IV}}t + 5\varepsilon^{\mathrm{V}} - 2\varepsilon^{\mathrm{IV}})
\end{aligned}
\]
and to $q^{\mathrm{V}}$ (ibid.)
\[
\begin{aligned}
&-127''13 - t.1''0212\\
&{}+ (179''952 - t.0''025192)\sin(5n^{\mathrm{V}}t - 2n^{\mathrm{IV}}t + 5\varepsilon^{\mathrm{V}} - 2\varepsilon^{\mathrm{IV}})\\
&{}- (116''541 + t.0''070196)\cos(5n^{\mathrm{V}}t - 2n^{\mathrm{IV}}t + 5\varepsilon^{\mathrm{V}} - 2\varepsilon^{\mathrm{IV}})
\end{aligned}
\]
This reduces the error in longitude mentioned above to $5'\,22''$, small enough to attribute to the observation without hesitation.

